\documentclass[12pt,a4paper]{report}

% Encodage
\usepackage[T1]{fontenc}
\usepackage[utf8]{inputenc}
\usepackage[french]{babel}

% Mise en page
\usepackage{geometry}
\usepackage{setspace}
\usepackage{lmodern}

% Couleurs et dessins
\usepackage{xcolor}
\usepackage{tikz}

% Tableaux
\usepackage{longtable}
\usepackage{booktabs}
\usepackage{array}

% Code source
\usepackage{listings}

% Images
\usepackage{graphicx}
\usepackage{float}

% Hyperliens
\usepackage{hyperref}

% Guillemets français
\usepackage{csquotes}

% Pour éviter les débordements
\usepackage{ragged2e}

\geometry{
    top=2.5cm,
    bottom=2.5cm,
    left=2.5cm,
    right=2.5cm
}

\setstretch{1.15}

%=================================================
% STYLE LISTINGS
%=================================================

\lstset{
    basicstyle=\ttfamily\small,
    breaklines=true,
    frame=single,
    columns=fullflexible
}

%=================================================
% PLACEHOLDERS DIAGRAMMES ET CAPTURES
%=================================================

\newcommand{\diagramplaceholder}[3]{
\begin{figure}[H]
\centering
\fbox{
\begin{minipage}[c][6cm][c]{0.8\textwidth}
\centering
\textbf{#2}

\vspace{0.5cm}

#3
\end{minipage}
}
\caption{#2}
\end{figure}
}

\newcommand{\screenshotplaceholder}[3]{
\begin{figure}[H]
\centering
\fbox{
\begin{minipage}[c][6cm][c]{0.8\textwidth}
\centering
\textbf{#2}

\vspace{0.5cm}

#3
\end{minipage}
}
\caption{#2}
\end{figure}
}

\begin{document}

%==================================
% TITRE PERSONNALISÉ
%==================================

\vspace*{0.5cm}

\noindent
\begin{tikzpicture}

% Rectangle noir
\fill[black] (0,0) rectangle (0.9,1.8);

% Ligne horizontale
\draw[line width=0.5pt] (1.2,0.9) -- (14.8,0.9);

% Titre
\node[
anchor=west,
font=\fontsize{24}{28}\bfseries
] at (1.2,0.25)
{TABLE DES MATIÈRES};

\end{tikzpicture}

\vspace{1.2cm}

\tableofcontents
\clearpage

%=================================================
% INTRODUCTION GENERALE
%=================================================

\chapter*{Introduction générale}
\addcontentsline{toc}{chapter}{Introduction générale}

Texte introduction...

%=================================================
% CHAPITRE 1
%=================================================

\chapter{Cadre général du projet}

\section{Introduction}

\section{Présentation de l'organisme d'accueil}

\section{Présentation du projet}
\subsection{Contexte du projet}
\subsection{Problématique}
\subsection{Étude critique de l'existant}
\subsection{Solution proposée}

\section{Méthodologie de travail}
\subsection{Méthodologie Agile Scrum}
\subsection{Planification globale du projet}

\section{Environnement technique}
\subsection{Technologies backend}
\subsection{Technologies frontend}
\subsection{Base de données et accès aux données}
\subsection{Outils de développement et de test}

\section{Conclusion}

%=================================================
% CHAPITRE 2
%=================================================

\chapter{Analyse et spécification des besoins}

\section{Introduction}

\section{Identification des acteurs}
\subsection{Super administrateur}
\subsection{Administrateur}
\subsection{Utilisateur final}

\section{Clarification du vocabulaire métier}
\subsection{Famille de documents}
\subsection{Modèle de document ou template}
\subsection{Charte graphique}
\subsection{Configuration de vue de table}
\subsection{Instance de document}

\section{Analyse des besoins fonctionnels}
\subsection{Gestion des utilisateurs}
\subsection{Gestion des rôles et permissions}
\subsection{Gestion des familles de documents}
\subsection{Gestion des modèles de documents}
\subsection{Gestion des chartes graphiques}
\subsection{Gestion des configurations de vues de tables}
\subsection{Édition et génération des documents}
\subsection{Gestion des bénéficiaires et des données métier}

\section{Analyse des besoins non fonctionnels}
\subsection{Sécurité}
\subsection{Performance}
\subsection{Ergonomie et expérience utilisateur}
\subsection{Maintenabilité}

\section{Backlog du produit}
\subsection{User stories principales}
\subsection{Priorisation des besoins}
\subsection{Critères d'acceptation}

\section{Planification des sprints}
\subsection{Découpage temporel}
\subsection{Objectifs par sprint}
\subsection{Livrables attendus}

\section{Diagrammes UML globaux}
\subsection{Diagramme de cas d'utilisation général}
\subsection{Diagramme de classes global}
\subsection{Diagramme de séquence global}

\section{Conclusion}

%=================================================
% CHAPITRE 3
%=================================================

\chapter{Sprint 1 : Authentification et socle applicatif}

\section{Introduction}

Le premier sprint constitue la fondation technique de la plateforme de gestion documentaire. Avant de traiter les modèles de documents, les chartes graphiques ou les bénéficiaires, il était indispensable de mettre en place un socle applicatif stable, sécurisé et capable de supporter les évolutions fonctionnelles prévues dans les sprints suivants. Ce sprint s'est donc concentré sur l'initialisation du backend ASP.NET Core Web API, la configuration de la base de données SQL Server, la mise en place de l'authentification JWT et la préparation des premiers points d'entrée de l'application.

Dans une application documentaire, l'authentification n'est pas une fonctionnalité périphérique. Elle conditionne l'accès aux modèles, aux données métier, aux documents générés et aux opérations d'administration. Une mauvaise gestion de l'identité peut entraîner une exposition de données sensibles ou un accès non autorisé aux documents internes. Pour cette raison, le sprint a été conçu autour d'un principe simple : chaque utilisateur doit pouvoir être identifié de manière fiable avant toute interaction avec le système.

Le résultat attendu de ce sprint est un backend opérationnel, documenté par Swagger, connecté à SQL Server et capable de générer des jetons JWT exploitables par l'application Angular. Cette base technique servira ensuite à appliquer la gouvernance d'accès par rôles dans le sprint suivant.

\section{Objectifs du sprint}

Les objectifs du sprint se répartissent en objectifs fonctionnels et objectifs techniques. Sur le plan fonctionnel, il s'agit de permettre à un utilisateur de se connecter, de récupérer son profil et d'être redirigé vers l'espace correspondant à son rôle. Sur le plan technique, il s'agit de structurer l'API, d'organiser les couches applicatives et de préparer les services communs utilisés par les modules métier.

Les objectifs retenus sont les suivants :

\begin{itemize}
    \item créer le projet backend ASP.NET Core Web API et organiser les dossiers \texttt{Controllers}, \texttt{Services}, \texttt{Repositories}, \texttt{DTOs}, \texttt{Domain} et \texttt{Infrastructure} ;
    \item configurer la connexion à SQL Server à travers une fabrique de connexion dédiée ;
    \item définir l'entité \texttt{User} ainsi que les objets de transfert associés à l'authentification ;
    \item implémenter l'inscription, la connexion et la récupération du profil utilisateur ;
    \item générer un jeton JWT contenant l'identifiant, le nom, l'adresse électronique, le rôle et éventuellement l'organisation de l'utilisateur ;
    \item protéger les premiers endpoints à l'aide du middleware d'authentification ;
    \item configurer CORS afin d'autoriser la communication entre l'interface Angular et l'API ;
    \item intégrer Swagger pour faciliter les tests et la documentation des endpoints.
\end{itemize}

Ces objectifs permettent d'obtenir une première version exécutable de l'application. Ils préparent également les éléments de sécurité nécessaires à la gestion des rôles et des autorisations.

\section{Backlog du sprint}

Le backlog du sprint regroupe les user stories prioritaires liées à la mise en place du socle applicatif. Chaque user story est associée à des tâches techniques permettant de transformer le besoin utilisateur en fonctionnalités vérifiables.

\begin{longtable}{p{0.17\textwidth} p{0.38\textwidth} p{0.35\textwidth}}
\caption{Backlog du sprint 1}\\
\toprule
\textbf{User story} & \textbf{Description} & \textbf{Tâches principales}\\
\midrule
\endfirsthead
\toprule
\textbf{User story} & \textbf{Description} & \textbf{Tâches principales}\\
\midrule
\endhead
US1.1 & En tant qu'utilisateur, je peux me connecter avec une adresse électronique et un mot de passe valides. & Créer \texttt{LoginRequest}, exposer \texttt{POST /api/auth/login}, vérifier le mot de passe, générer le JWT, retourner les informations du profil.\\
\midrule
US1.2 & En tant que nouvel utilisateur, je peux créer un compte afin d'accéder à la plateforme. & Créer \texttt{RegisterRequest}, vérifier l'unicité de l'utilisateur, hacher le mot de passe avec BCrypt, enregistrer l'utilisateur en base.\\
\midrule
US1.3 & En tant qu'utilisateur authentifié, je peux consulter mon profil. & Protéger \texttt{GET /api/auth/profile}, lire l'identifiant depuis les claims JWT, retourner un \texttt{UserResponse}.\\
\midrule
US1.4 & En tant que développeur, je peux tester les endpoints à partir d'une documentation interactive. & Configurer Swagger, ajouter la définition de sécurité Bearer, tester les appels authentifiés.\\
\midrule
US1.5 & En tant qu'application Angular, je peux communiquer avec l'API de manière sécurisée. & Configurer CORS, mettre en place l'interceptor Angular, stocker le token côté client et l'injecter dans les requêtes.\\
\bottomrule
\end{longtable}

\subsection{Authentification JWT}

La user story d'authentification JWT représente le coeur de ce sprint. L'utilisateur fournit ses identifiants à travers le formulaire de connexion. Le backend vérifie l'existence du compte, compare le mot de passe reçu avec le mot de passe haché en base, puis génère un jeton signé. Ce jeton devient la preuve d'authentification utilisée dans les appels suivants.

Le JWT transporte des claims essentiels : l'identifiant technique de l'utilisateur, son nom, son rôle, son adresse électronique et son organisation si elle existe. Ces informations évitent d'interroger la base à chaque requête pour les contrôles simples, tout en laissant au backend la possibilité de vérifier les informations sensibles lorsque cela est nécessaire.

\subsection{Gestion du profil utilisateur}

La gestion du profil consiste à exposer un endpoint permettant de récupérer les informations de l'utilisateur connecté. Cette fonctionnalité est utile pour l'affichage de l'espace personnel, la redirection selon le rôle et la personnalisation des pages Angular. Elle permet également de vérifier que le jeton envoyé par le frontend est bien valide.

\subsection{Configuration du backend}

Le backend est structuré selon une séparation claire des responsabilités. Les contrôleurs reçoivent les requêtes HTTP, les services portent la logique applicative, les repositories gèrent l'accès aux données, et les DTOs définissent les contrats d'entrée et de sortie. Cette organisation prépare l'application à accueillir progressivement les modules métier des sprints suivants.

\subsection{Connexion à la base de données}

La connexion SQL Server est centralisée dans une fabrique dédiée. Cette approche évite de disperser les chaînes de connexion dans le code et facilite la maintenance. Elle prépare également l'utilisation de Dapper pour exécuter les requêtes SQL de manière performante et contrôlée.

\section{Analyse et conception}

L'analyse de ce sprint porte sur le cycle d'authentification et sur les composants nécessaires au démarrage sécurisé de l'application. L'objectif est de concevoir une chaîne complète allant de la saisie des identifiants dans Angular jusqu'à la validation du token par ASP.NET Core.

\subsection{Diagramme de cas d'utilisation}

\diagramplaceholder{sprint1_auth_socle}{Diagramme de cas d'utilisation du sprint 1}{Cas d'utilisation liés à l'authentification et au socle applicatif}

Le diagramme de cas d'utilisation du sprint 1 présente deux acteurs principaux : l'utilisateur non authentifié et l'utilisateur authentifié. L'utilisateur non authentifié peut s'inscrire, se connecter et consulter la documentation Swagger dans le cadre des tests techniques. Après authentification, l'utilisateur peut consulter son profil et accéder aux fonctionnalités protégées.

Le cas d'utilisation \og Se connecter \fg{} inclut la vérification des identifiants et la génération du jeton JWT. Le cas d'utilisation \og Consulter le profil \fg{} nécessite une authentification préalable, car le backend doit extraire l'identité de l'utilisateur depuis les claims du token.

Cette modélisation montre que le sprint ne se limite pas à un simple formulaire de connexion. Il définit le point d'entrée sécurisé de toute la plateforme et conditionne l'accès à l'ensemble des modules à venir.

\subsection{Diagramme de classes}

\diagramplaceholder{sprint1_auth_socle}{Diagramme de classes du sprint 1}{Classes principales du socle d'authentification}

Le diagramme de classes met en évidence la collaboration entre \texttt{AuthController}, \texttt{AuthService}, \texttt{IUserRepository}, \texttt{User}, \texttt{LoginRequest}, \texttt{RegisterRequest}, \texttt{AuthResponse} et \texttt{UserResponse}. Le contrôleur dépend d'une abstraction de service, ce qui réduit le couplage entre la couche HTTP et la logique métier.

L'entité \texttt{User} contient les informations persistées : identifiant, nom d'utilisateur, adresse électronique, mot de passe haché, rôle, organisation, profil, date de création et état du compte. Les DTOs évitent d'exposer directement l'entité de domaine dans les réponses HTTP. Ainsi, le champ \texttt{PasswordHash} reste interne au backend et n'est jamais transmis au frontend.

L'utilisation de l'interface \texttt{IAuthService} permet également d'isoler la logique d'authentification. Cette séparation facilite les tests unitaires et rend possible une évolution future vers un autre mécanisme d'identité sans modifier les contrôleurs.

\subsection{Diagrammes de séquence}

\diagramplaceholder{sprint1_auth_socle}{Diagramme de séquence du sprint 1}{Séquence de connexion et récupération du profil}

Le diagramme de séquence décrit le scénario principal de connexion. L'utilisateur saisit ses identifiants dans le composant Angular de connexion. Le composant appelle \texttt{AuthService}, qui transmet la requête à \texttt{ApiService}. L'API reçoit la requête dans \texttt{AuthController}, puis délègue le traitement à \texttt{AuthService} côté backend. Ce service consulte \texttt{UserRepository}, vérifie le mot de passe et génère le token.

Lorsque la réponse revient au frontend, le jeton et les informations utilisateur sont stockés localement. Les requêtes suivantes passent par un interceptor HTTP qui ajoute l'en-tête \texttt{Authorization: Bearer}. Lorsqu'un endpoint protégé est appelé, le middleware ASP.NET Core valide automatiquement la signature, l'émetteur, l'audience et la durée de validité du token.

\subsubsection*{Description des interactions système}

Le système met en jeu deux chaînes d'interaction. La première est synchrone et concerne l'authentification : formulaire, service Angular, API, repository et base de données. La seconde est transversale : elle concerne tous les appels ultérieurs et passe par l'interceptor Angular et le middleware d'authentification ASP.NET Core.

Cette conception réduit la duplication. Le frontend n'a pas à ajouter manuellement le token dans chaque service métier. De la même manière, chaque contrôleur backend n'a pas à réimplémenter la validation du token, car cette responsabilité est prise en charge par le middleware.

\subsubsection*{Analyse technique}

La sécurité repose sur une signature symétrique HMAC SHA-256 configurée dans \texttt{JwtSettings}. Le paramètre \texttt{ClockSkew} est fixé à zéro afin d'éviter une tolérance excessive après expiration du jeton. Les mots de passe sont hachés avec BCrypt lors de l'inscription, ce qui permet de stocker une représentation non réversible du secret utilisateur.

Le backend utilise l'injection de dépendances d'ASP.NET Core pour enregistrer les services et repositories. Cette approche permet de créer des instances par requête et d'éviter le couplage direct entre les composants.

\subsubsection*{Architecture des composants concernés}

L'architecture du sprint peut être résumée en quatre niveaux. Le niveau présentation est représenté par le composant Angular de connexion. Le niveau communication est assuré par \texttt{ApiService} et l'interceptor HTTP. Le niveau applicatif backend est constitué par \texttt{AuthController} et \texttt{AuthService}. Enfin, le niveau persistance regroupe \texttt{UserRepository}, la fabrique de connexion et SQL Server.

\section{Réalisation}

La réalisation a été menée en commençant par le backend afin de disposer d'une API testable avec Swagger, puis en intégrant progressivement l'interface Angular. Cette démarche a permis de valider les contrats HTTP avant de connecter le frontend.

\subsection{Mise en place du backend ASP.NET Core}

Le backend a été initialisé sous forme d'une Web API ASP.NET Core. Le fichier \texttt{Program.cs} centralise la configuration des services, du CORS, de l'authentification, de Swagger et de l'injection de dépendances. Les contrôleurs sont activés par \texttt{AddControllers()}, tandis que les services métier et repositories sont enregistrés en portée \texttt{Scoped}.

\begin{lstlisting}[caption={Configuration JWT et injection des services dans ASP.NET Core}]
builder.Services.AddAuthentication(options =>
{
    options.DefaultAuthenticateScheme = JwtBearerDefaults.AuthenticationScheme;
    options.DefaultChallengeScheme = JwtBearerDefaults.AuthenticationScheme;
})
.AddJwtBearer(options =>
{
    options.TokenValidationParameters = new TokenValidationParameters
    {
        ValidateIssuer = true,
        ValidateAudience = true,
        ValidateLifetime = true,
        ValidateIssuerSigningKey = true,
        ValidIssuer = jwtSettings?.Issuer,
        ValidAudience = jwtSettings?.Audience,
        IssuerSigningKey = new SymmetricSecurityKey(
            Encoding.UTF8.GetBytes(jwtSettings?.SecretKey ?? "")),
        ClockSkew = TimeSpan.Zero
    };
});
\end{lstlisting}

Cette configuration impose la validation de l'émetteur, de l'audience, de la durée de vie et de la clé de signature. Elle garantit que seuls les jetons générés par l'application et encore valides permettent d'accéder aux endpoints protégés.

\subsection{Implémentation de l'authentification}

L'authentification est exposée à travers \texttt{AuthController}. Le endpoint \texttt{POST /api/auth/login} reçoit un \texttt{LoginRequest}, délègue la vérification à \texttt{AuthService}, puis retourne un \texttt{AuthResponse}. Le endpoint \texttt{GET /api/auth/profile} est protégé par \texttt{[Authorize]} et récupère l'identifiant utilisateur à partir des claims.

\begin{lstlisting}[caption={Endpoint de récupération du profil utilisateur}]
[HttpGet("profile")]
[Authorize]
public async Task<ActionResult<UserResponse>> GetProfile()
{
    var userIdClaim = User.FindFirst(ClaimTypes.NameIdentifier)?.Value;
    if (string.IsNullOrEmpty(userIdClaim) ||
        !int.TryParse(userIdClaim, out var userId))
    {
        return Unauthorized();
    }

    var profile = await _authService.GetUserProfileAsync(userId);
    return Ok(profile);
}
\end{lstlisting}

La génération du token est centralisée dans le service d'authentification. Les claims insérés permettent au backend et au frontend de connaître le rôle et l'organisation de l'utilisateur sans exposer de données sensibles.

\begin{lstlisting}[caption={Claims principaux intégrés au jeton JWT}]
var claims = new List<Claim>
{
    new(ClaimTypes.NameIdentifier, userId.ToString()),
    new(ClaimTypes.Name, username),
    new(ClaimTypes.Role, role),
    new(JwtRegisteredClaimNames.Jti, Guid.NewGuid().ToString()),
    new(JwtRegisteredClaimNames.Iat,
        DateTimeOffset.UtcNow.ToUnixTimeSeconds().ToString(),
        ClaimValueTypes.Integer64)
};
\end{lstlisting}

Côté Angular, \texttt{AuthService} encapsule la connexion, la déconnexion, le stockage du token et la récupération de l'utilisateur courant. Le token est stocké dans le \texttt{localStorage}, puis lu par l'interceptor HTTP.

\begin{lstlisting}[caption={Ajout automatique du token dans les requêtes Angular}]
export const authInterceptor: HttpInterceptorFn = (req, next) => {
  const authService = inject(AuthService);
  const token = authService.getToken();

  if (token && !req.url.includes('/auth')) {
    const authReq = req.clone({
      headers: req.headers.set('Authorization', `Bearer ${token}`)
    });
    return next(authReq);
  }

  return next(req);
};
\end{lstlisting}

\screenshotplaceholder{sprint1_login.png}{Interface de connexion}{Interface de connexion de la plateforme}
\screenshotplaceholder{sprint1_profile.png}{Profil utilisateur authentifié}{Consultation du profil après authentification}

\subsection{Configuration CORS, Swagger et middlewares}

La politique CORS nommée \texttt{EditorCors} permet au frontend Angular de communiquer avec l'API. Elle autorise les méthodes et en-têtes nécessaires aux appels REST et active la transmission des informations d'authentification.

Swagger est configuré avec une définition de sécurité Bearer. Cette configuration permet de tester les endpoints protégés directement depuis l'interface Swagger UI en renseignant le jeton JWT obtenu après connexion.

L'ordre des middlewares est également important. La redirection HTTPS est appliquée avant CORS, puis viennent l'authentification et l'autorisation. Les contrôleurs sont ensuite mappés par \texttt{MapControllers()}.

\subsection{Configuration de la base SQL Server}

La persistance repose sur SQL Server. Le backend utilise une fabrique de connexion afin de créer les connexions au moment de l'exécution des requêtes. Cette fabrique permet d'isoler les paramètres de connexion dans la configuration et de faciliter les changements d'environnement.

Les données d'authentification sont stockées dans la table des utilisateurs. Le mot de passe n'est jamais stocké en clair. La colonne \texttt{PasswordHash} contient le résultat du hachage BCrypt. Les champs \texttt{Role}, \texttt{OrganizationId}, \texttt{Profile} et \texttt{ProfileDetail} préparent la personnalisation des accès dans les sprints suivants.

\section{Tests et validation}

Les tests du sprint ont porté sur les scénarios fonctionnels, la sécurité et l'intégration Angular/API.

\begin{longtable}{p{0.24\textwidth} p{0.42\textwidth} p{0.24\textwidth}}
\caption{Scénarios de test du sprint 1}\\
\toprule
\textbf{Scénario} & \textbf{Description} & \textbf{Résultat attendu}\\
\midrule
\endfirsthead
\toprule
\textbf{Scénario} & \textbf{Description} & \textbf{Résultat attendu}\\
\midrule
\endhead
Connexion valide & Saisir une adresse électronique et un mot de passe corrects. & Token retourné, utilisateur stocké, redirection correcte.\\
\midrule
Connexion invalide & Saisir un mauvais mot de passe. & Réponse d'erreur, aucun token stocké.\\
\midrule
Profil sans token & Appeler \texttt{/api/auth/profile} sans en-tête Authorization. & Réponse \texttt{401 Unauthorized}.\\
\midrule
Profil avec token & Appeler \texttt{/api/auth/profile} avec un token valide. & Informations du profil retournées.\\
\midrule
Expiration du token & Utiliser un token expiré. & Rejet de la requête par le middleware JWT.\\
\bottomrule
\end{longtable}

Les validations techniques ont confirmé que le token était ajouté automatiquement par l'interceptor Angular, que Swagger acceptait l'authentification Bearer et que les endpoints protégés refusaient les appels anonymes. Les tests de sécurité ont également vérifié que le mot de passe haché n'était jamais présent dans les réponses HTTP.

\section{Conclusion}

Ce premier sprint a permis de mettre en place le socle technique de la plateforme. L'application dispose désormais d'une API ASP.NET Core structurée, d'une connexion SQL Server, d'un système d'authentification JWT, d'une documentation Swagger et d'une intégration Angular capable de stocker et transmettre le jeton d'accès.

Cette base est essentielle pour la suite du projet. Elle permet d'aborder le sprint suivant dans de bonnes conditions, en construisant la gouvernance d'accès sur une identité utilisateur déjà fiable. Le prochain sprint sera donc consacré à la gestion des rôles, à la protection des routes Angular et à la sécurisation des endpoints selon les profils \texttt{supAdmin}, \texttt{admin} et \texttt{user}.



%=================================================
% CHAPITRE 4
%=================================================

\chapter{Sprint 2 : Gestion des rôles et gouvernance d'accès}

\section{Introduction}
\section{Objectifs du sprint}
\section{Backlog du sprint}
\subsection{Définition des rôles}
\subsection{Protection des routes et des endpoints}
\subsection{Gestion des utilisateurs et profils}
\section{Analyse et conception}
\subsection{Diagramme de cas d'utilisation}
\subsection{Diagramme de classes}
\subsection{Diagrammes de séquence}
\section{Réalisation}
\subsection{Implémentation du système RBAC}
\subsection{Mise en place des autorisations backend}
\subsection{Gestion des guards côté frontend}
\subsection{Sécurisation des accès par profil}
\section{Tests et validation}
\section{Conclusion}

%=================================================
% CHAPITRE 5
%=================================================

\chapter{Sprint 3 : Module coeur d'édition documentaire}

\section{Introduction}
\section{Objectifs du sprint}
\section{Backlog du sprint}
\subsection{Intégration de l'éditeur riche}
\subsection{Création et modification du contenu documentaire}
\subsection{Sauvegarde et chargement des contenus}
\subsection{Prévisualisation du document}
\section{Analyse et conception}
\subsection{Diagramme de cas d'utilisation}
\subsection{Diagramme de classes}
\subsection{Diagrammes de séquence}
\section{Réalisation}
\subsection{Intégration de l'éditeur Tiptap}
\subsection{Gestion du contenu des modèles de documents}
\subsection{Application des en-têtes, corps et pieds de page}
\subsection{Prévisualisation et rendu des documents}
\section{Tests et validation}
\section{Conclusion}

%=================================================
% CHAPITRE 6
%=================================================

\chapter{Sprint 4 : Administration des modules métier}

\section{Introduction}
\section{Objectifs du sprint}
\section{Backlog du sprint}
\subsection{Gestion des familles de documents}
\subsection{Gestion des modèles de documents}
\subsection{Gestion des chartes graphiques}
\subsection{Gestion des configurations de vues de tables}
\subsection{Gestion des bénéficiaires et filtres métier}
\section{Analyse et conception}
\subsection{Diagramme de cas d'utilisation}
\subsection{Diagramme de classes}
\subsection{Diagrammes de séquence}
\section{Réalisation}
\subsection{CRUD des familles de documents}
\subsection{CRUD des modèles de documents}
\subsection{CRUD des chartes graphiques}
\subsection{CRUD des configurations de vues de tables}
\subsection{Association famille, modèle et charte graphique}
\subsection{Configuration des données métier et des filtres}
\section{Tests et validation}
\section{Conclusion}

%=================================================
% CHAPITRE 7
%=================================================

\chapter{Sprint 5 : Intégration, qualité et livraison}

\section{Introduction}
\section{Objectifs du sprint}
\section{Backlog du sprint}
\subsection{Intégration frontend/backend}
\subsection{Tests fonctionnels et tests d'intégration}
\subsection{Optimisation des performances}
\subsection{Préparation de la livraison}
\section{Réalisation}
\subsection{Validation des parcours utilisateurs}
\subsection{Gestion des erreurs}
\subsection{Optimisation et renforcement de la sécurité}
\subsection{Documentation technique et utilisateur}
\subsection{Préparation de la démonstration PFE}
\section{Tests et validation}
\section{Conclusion}

%=================================================
% CONCLUSION
%=================================================

\chapter*{Conclusion générale}
\addcontentsline{toc}{chapter}{Conclusion générale}

Texte conclusion...

%=================================================
% WEBOGRAPHIE
%=================================================

\chapter*{Webographie}
\addcontentsline{toc}{chapter}{Webographie}

%=================================================
% BIBLIOGRAPHIE
%=================================================

\chapter*{Bibliographie}
\addcontentsline{toc}{chapter}{Bibliographie}

%=================================================
% ANNEXES
%=================================================

\chapter*{Annexes}
\addcontentsline{toc}{chapter}{Annexes}

\end{document}
